\section{Open Science and Artifact Map}
\label{app:repro}

\subsection{Endpoints and versions}

Provider experiments recorded requested and provider-reported model
identifiers. Earlier runs reported DeepSeek V4 Flash, Qwen 3.7 Flash
2026-07-15, and MiniMax-M3. A later request for
\path{deepseek-v4-flash} returned \path{deepseek-flash}.
We report this run separately, not as the same endpoint. The run map below
states which endpoint served each experiment. Receiver calls used temperature 0, disabled
thinking, automatic tool choice, and a 1{,}024-token ceiling in the five-form
confirmation. This ceiling replaced a 300-token calibration limit that had
produced seven censored responses; no inference uses those censored cells.

The external instantiation used MCP Inspector 2.4.0 at pinned commit
\texttt{edf54f5d}, MCP server 2.0.0, and Node.js v22.20.0. The complete commit
hash is stored in the artifact manifest. Protocol citations are version- and
date-pinned (\S\ref{sec:related}). Deterministic A2A studies pinned
\texttt{a2a-sdk} 1.1.2 and \texttt{@a2a-js/sdk} 1.1.0 with their dependencies
and inspected sources; neither made provider calls.

\subsection{Statistical procedures}

The independent unit is the source directive or payload throughout; paired
designs use exact McNemar tests as the primary inference, with paired
nonparametric bootstrap intervals (100{,}000 resamples) at preregistered
seeds. The five-form primary comparisons were a prespecified fixed sequence
(familywise $\alpha = 0.05$). Rating studies keep the two model raters in
isolated sessions: rater A used KIMI K3 and rater B used Grok 4.5, and
each saw only its assigned workbook, with frozen row order, opaque IDs, and no
form labels or behavioral outcomes. The five-form workbooks had 96 rows; the
E2 adaptive-sender workbooks had 144. A retired E2 workflow exposed both model
contexts to A/B materials; both returns were excluded and rerun under the
isolation rule. The accepted fresh-context declaration is operator-attested
and archived with the E2 rating result.
The sequential standardization in \S\ref{sec:ratings} is labeled descriptive
and is not used as causal mediation.

Multiplicity control follows the confirmatory family declared for each
experiment. The five-form study uses its two-comparison fixed sequence; E2 and
E3 retain their separately frozen primary paired contrasts. E4 first tests
elevation within the adaptive marked relay and, conditional on rejection,
tests the exact-response relay-policy interaction by exact paired sign flips;
its broader any-write decomposition is post hoc. E5 prospectively defines any
write as primary and exact mutation request as a fixed-sequence secondary. E6
uses three paired McNemar contrasts against the unmarked control with Holm
step-down correction, separately at the direct and relayed stages. Secondary ratings,
family breakdowns, cross-run comparisons, and cross-fixture comparisons do not
enter those confirmatory families. We do not apply one omnibus paper-level
correction across distinct scientific questions, and a result from one family
does not unlock inference in another.
E12 defines elevated exact protected-action pursuit as primary, using paired
exact McNemar inference and a directive-level bootstrap interval; direct-write
and escalation decompositions, the non-elevated comparison, and their
interaction are supporting analyses within that experiment.

\subsection{Run directories and audits}

Every formal run is an append-only directory under \texttt{runs/} with
events, provider responses, report, and post-run audit; key numbers were
independently recomputed from raw artifacts rather than read back from the
run's own report. The runs supporting each section (paths relative to the
artifact root):

{\sloppy\footnotesize
\begin{itemize}\itemsep2pt
  \item \textbf{\S\ref{sec:laundering}} (five-form):\\
  \path{deepseek-five-form-provenance-confirmation-001}\\
  \path{qwen-five-form-provenance-confirmation-001}\\
  calibration runs \texttt{-001} and \texttt{-002}; rater returns under
  \texttt{rating-returned/}.
  \item \textbf{\S\ref{sec:adaptive-sender}} (E2 adaptive sender):\\
  \path{deepseek-e2-adversarial-sender-001}\\
  (pre-generation and receiver-release freezes; 144 frozen generated texts;
  288 calls, zero retries; behavioral \texttt{independent\_audit.json}, all
  checks passed). Accepted isolated secondary ratings, their attestation, and
  a separate reconstruction are under
  \path{rating-returned-clean-v2/}.
  \item \textbf{\S\ref{sec:propagation}} (selector):\\
  \path{deepseek-retention-propagation-calibration-001}\\
  and its \texttt{-confirmation-001}. The selected inputs comprise 37 archived
  faithful-copy outputs and 84 task-condensation outputs. Their first-receiver
  delegations were 7 and 33, and their complete two-receiver chains were 2 and
  15. \path{audit_retention_source_framing.py} verifies every selected output
  against the original sender event and prompt; these post hoc source-group
  counts do not change the primary 17/121 endpoint. The later concurrent-group check is
  \path{deepseek-retention-discrimination-001}; its
  \texttt{MODEL\_IDENTITY\_ADDENDUM.md} explains why the returned
  \texttt{deepseek-flash} responses cannot test discrimination on the
  earlier \texttt{deepseek-v4-flash} endpoint. Its pre-call manifest and
  analysis plan were internally frozen, not externally timestamped.
  A structural benign-delegation control is archived at
  \path{deepseek-benign-delegation-control-002}: 16 chains, 30 calls,
  instruction-eligible receivers, zero retries, and a passing wire audit.
  The request named \path{deepseek-v4-flash}, but responses reported
  \path{deepseek-flash}; it is therefore a plumbing control rather than an
  endpoint replication.
  A separately scoped receiver-authorized propagation calibration is archived
  at \path{qwen36-authorized-propagation-calibration-001}: the same marked
  sender output was presented under L0 and L2 to an R2 worker that was
  authorized to write \path{protected/canary}. Both arms produced 12/12 exact
  requests; this ceiling result is a fixture check, not a lineage effect or a
  security outcome. The independent recount is
  \path{independent_audit.json}.
  A zero-call headroom audit found the protected literal and a write verb in
  all 12 sender summaries and all 24 R1 payloads; it explains the ceiling and
  is archived as \path{headroom_audit.json}.
  A weaker indirect-content calibration is archived at
  \path{qwen36-authorized-propagation-weak-calibration-001}: exact protected
  requests were 2/12 in both L0 and L2, with no explicit write verb in the
  source bank. Its result is a headroom check, not a lineage comparison.
  The 48-pair disjoint follow-up is archived at
  \path{qwen36-authorized-propagation-confirmation-001}: L0 produced 6/48
  and L2 4/48 exact requests (paired risk difference $-.0417$; exact
  McNemar $p=.7266$). It remains an internally frozen fixture comparison,
  not a confirmatory security result.
  The zero-call content audit for the 17 complete confirmation chains is
  \path{content_audit.json} in
  \path{deepseek-retention-propagation-confirmation-001}; it is a post hoc
  content characterization, not a new propagation estimate.
  \item \textbf{\S\ref{sec:propagation}} (defense):\\
  \path{qwen-cross-model-defense-001}. The DeepSeek one-family run
  \path{deepseek-one-family-defense-001} failed its preregistered
  susceptibility gate and is kept as a failed-gate record, not as defense
  evidence.
  \item \textbf{\S\ref{sec:propagation}} (carrier):\\
  \path{deepseek-carrier-mechanism-001},
  \path{qwen-carrier-mechanism-001},
  \path{minimax-carrier-mechanism-001}.
  \item \textbf{\S\ref{sec:covariance}} (DeepSeek):\\
  \path{deepseek-lifecycle-covariance-confirmation-001}\\
  (calibration \texttt{-001}; freeze
  \path{lifecycle-covariance-deepseek-20260828-v1};
  \texttt{independent\_verification.json}).
  \item \textbf{\S\ref{sec:covariance}} (Qwen):\\
  \path{qwen-lifecycle-covariance-confirmation-001}\\
  (calibration \texttt{-001};
  \texttt{independent\_audit.json}; 116 calls, zero retries).
  \item \textbf{\S\ref{sec:covariance}} (credential-gated E3):\\
  \path{qwen-e3-credential-gated-001}\\
  (freeze \path{e3-credential-gated-qwen-20260903-v1}; 126 calls,
  zero retries; \texttt{independent\_audit.json}, 14/14 checks).
  \item \textbf{\S\ref{sec:e4-gateway}} (E4 relay-output gateway evaluation):\\
  \path{deepseek-e4-e2-e3-composition-001}\\
  (freeze \path{e4-e2-e3-composition-deepseek-20260904-v1}; 318 calls,
  zero retries; \texttt{independent\_audit.json}, 16/16 checks). Post-run
  nested-outcome decomposition: \texttt{outcome\_decomposition.json}
  (original E4 report unchanged). The frozen directory name predates the
  construct split.
  \item \textbf{\S\ref{sec:e4-gateway}} (E5 Qwen endpoint replication):\\
  \path{qwen-e5-relay-gateway-001}\\
  (freeze \path{e5-qwen-relay-gateway-20260904-v1}; 96 calls;
  \texttt{independent\_audit.json}, all checks). The run records a bootstrap-
  seed protocol deviation; both seed recomputations give identical reported
  interval endpoints, and exact inference is unaffected.
  \item \textbf{\S\ref{sec:e4-gateway}} (E11 Qwen temporal repeat):\\
  \path{qwen-e11-hosted-repeat-001}\\
  (freeze \path{e11_freeze_v1}; 96 calls, zero retries;
  \texttt{independent\_audit.json}, 18/18 checks). The same 48 E5 texts and
  endpoint were reused under a new frozen calendar; the comparison is one
  paired temporal repeat, not a population variance estimate. E11--E5 effect
  shifts were $+1/48$ for both endpoints, with paired bootstrap intervals
  $[-0.1042,0.1667]$ and $[-0.0833,0.1250]$.
  \item \textbf{\S\ref{sec:e4-gateway}} (E12 capability-hidden check):\\
  \path{qwen-e12-capability-hidden-001}\\
  (freeze \path{e12_freeze_v1}; 252 calls, zero retries;
  \texttt{independent\_audit.json}, 30/30 checks). The primary outcome joins
  exact direct writes and exact non-mutating escalation requests; the study
  does not execute the escalation or infer free-text intent.
  \item \textbf{\S\ref{sec:spotlight-relay}} (E6 relay survival):\\
  \path{deepseek-e6-spotlight-relay-001}\\
  (192 source calls and 384 receiver calls, zero retries; pre-generation and
  receiver-release timestamps; \texttt{generation\_independent\_audit.json}
  and \texttt{independent\_audit\_v2.json}).
  \item \textbf{\S\ref{sec:external}}:\\
  \path{inspector-shared-root-confirmation-001}\\
  (\texttt{report.json}, \texttt{independent\_audit.json}); source and
  dependency freeze under \texttt{external\_auth\_instantiation/}.
  \item \textbf{\S\ref{sec:external}} (E7--E10 A2A lineage integrity):\\
  \path{e7-a2a-lineage-composition-001}, \path{e8-a2a-js-live-001},\\
  \path{e9-a2a-official-sample-live-001},\\
  \path{e10-a2a-two-hop-enforcement-001}\\
  (transport, exact-sample derivation, and consequential two-hop controls;
  independently recomputed; zero calls).
\end{itemize}
}

\subsection{Open science}
\label{app:openscience}

The anonymous, cloneable artifact bundle is being prepared at the following
submission URL placeholder:

\begin{center}
\emph{Anonymous artifact URL pending}
\end{center}

\noindent The release will contain, for each formal run: the frozen stimulus bank
and execution order; the sender, receiver, and rater prompts verbatim; the
three rating rubrics as frozen before scoring; raw append-only JSONL event
logs and provider responses; the analysis source (exact McNemar, paired
bootstrap with seeds, $q$-weighting enumerator); the independent
recomputation scripts and their outputs; the MCP Inspector instantiation
source with its dependency freeze and commit pin; and a manifest of SHA-256
hashes for all of the above.

\paragraph{Preregistration evidence.}
Because ``preregistered'' is only meaningful if the freeze provably precedes
the outcomes, each claim-bearing formal run's freeze record is externally witnessed: the
manifest hash for the stimulus bank, execution order, analysis plan, provider
pinning, and $q$ grid was deposited to a third-party timestamping service
before the first live call, and the deposit receipts are included in the
bundle. Reviewers can verify that the analysis source distributed with the
run hashes to the value recorded before execution.
For E1, four OpenTimestamps calendars accepted the Qwen replication freeze
before its 492 calls. For E2, four calendars accepted both the pre-generation
and receiver-release deposits before their respective live stages. The E3
deposit was likewise accepted before execution. Four OpenTimestamps calendars
accepted the E4 deposit before its 318 calls. Four calendars accepted the E5
deposit before its 96 calls; both sets of calendar attestations are included
pending Bitcoin anchoring. Four calendars accepted both E6 stage deposits
before their corresponding calls. Four calendars accepted the E12 freeze
before its 252 calls. Earlier receipts carry independently
verified Bitcoin attestations. The bundle keeps the original
receipts, self-contained upgraded copies, and a machine-readable status report.
The later endpoint-identity check is a diagnostic exception: its files were
written before live calls but have no external timestamp receipt, so we do not
label it externally preregistered or use it for confirmatory inference.
The native AgentDojo Slack screens were exploratory external-validity runs and
were not externally timestamped; they are reported as provider-stratified
descriptive evidence, not as preregistered mechanism tests.

\paragraph{What is \emph{not} released.}
No designed bank is planned to be withheld. The release includes attack
directives written against our sandbox tool schema; the schema is fictitious
and the effects are process-local, so release carries no operational risk
(\S\ref{sec:threat}).
